\begin{proof}
Clause numbers refer to the formal statement: (i) Step 1, (ii) Steps 2 to 5, (iii) Step 6, (iv) Step 9, (v) Step 10, (vi) Step 7, (vii) Step 8. In index reading $T(f)$ is $(\partial_{\mu}\partial_{\nu} - g_{\mu\nu}\partial^{2})f$, and the displayed formulas use the index reading.

Throughout, $s = Q(x)$, $L_{\mu} = \ln(\mu^{2}s)$, and primes are derivatives in $s$. Off the origin, for a function $f(s)$, $\partial_{\mu}s = 2x_{\mu}$ and $\partial^{\mu}x_{\mu} = 4$ give, in either signature,
\[
\partial_{\mu}\partial_{\nu}f = 2f'g_{\mu\nu} + 4f''x_{\mu}x_{\nu}, \qquad \partial^{2}f = 8f' + 4sf'', \qquad \partial^{\mu}(Ag_{\mu\nu} + Bx_{\mu}x_{\nu}) = (2A' + 2sB' + 5B)x_{\nu}.
\]
Every identity below between kernels off the origin is an identity between finite sums $\sum c\,s^{a}L_{\mu}^{k}$ (times $g$, $xx$ or $x$) and is computed exactly in the check.

\emph{Step 1 (traces and degrees).} In $\Cl(1,3)$ and in $\Cl(4,0)$, $\gamma_{\mu}x\gamma_{\nu}x = \gamma_{\mu}(2x_{\nu} - \gamma_{\nu}x)x = 2x_{\nu}\gamma_{\mu}x - \gamma_{\mu}\gamma_{\nu}Q(x)$, whose scalar part is $2x_{\mu}x_{\nu} - g_{\mu\nu}Q(x)$; a product of three vectors has no scalar part; $\gamma^{\mu}x\gamma_{\mu} = -2x$ in four dimensions. Left multiplication by a basis blade other than $1$ permutes the blades with signs and fixes none, so its real trace on the algebra is $16\langle X\rangle_{0}$; the algebra is the sum of two copies of the Dirac module, which has real dimension $8$ and complex dimension $4$, so the matrix trace on it is $4\langle X\rangle_{0}$. The massless Dirac propagator of Proposition \claimref{field-theory.second-quantisation} on the Euclidean section is $S_{0} = -\partial\!\!\!/\,(1/(4\pi^{2}s)) = x/(2\pi^{2}s^{2})$, since $-\partial^{2}(1/(4\pi^{2}s)) = \delta$ (Step 3); it is homogeneous of degree $-3$, so $\mathrm{sd}(S_{0}) = 3$, and the scalar propagator has $\mathrm{sd} = 2$. A product of kernels in one relative variable has scaling degree the sum: $6$ for $S_{0}\otimes S_{0}$, $5$ for $S_{0}D_{0}$, $4$ for $D_{0}^{2}$; with $n$ the number of relative coordinates times four, $\omega = \mathrm{sd} - n$ is $2, 1, 0$, and power counting gives $0$, $1$, $0$ for the vertex ($3 + 3 + 2 - 8$), the three-current ($9 - 8$) and the four-photon ($12 - 12$) kernels.

\emph{Step 2 (the massive propagator to order $m^{2}$).} On the flat field the Hadamard form (Fact \claimref{synthesis.imported-hadamard-parametrix-of-the-squared-dirac-operator-leads-with-van-vleck-root-times-spin-transport}) gives the scalar kernel $\Delta = \frac{1}{4\pi^{2}s} + V\ln s + W$ with $V$, $W$ smooth and $V(0) = Am^{2}$; the Klein-Gordon equation $(-\partial^{2} + m^{2})\Delta = 0$ off the origin at order $m^{2}$ reads $-A\,\partial^{2}\ln s + \frac{1}{4\pi^{2}s} = 0$, and $\partial^{2}\ln s = 4/s$, so $A = 1/(16\pi^{2})$. With $S = (-\partial\!\!\!/ + m)\Delta$ (so that $(\partial\!\!\!/ + m)S = \delta$), $S = \alpha x + \beta + R$ with $\alpha = \frac{1}{2\pi^{2}s^{2}} - \frac{m^{2}}{8\pi^{2}s}$, $\beta = \frac{m}{4\pi^{2}s}$, and $R$ the rest: its vector part is $x$ times $O(m^{4}\ln s)$ and its scalar part $O(m^{3}\ln s)$, both of scaling degree at most $0$ (a power of $\ln s$ does not raise the scaling degree).

\emph{Step 3 (the delta normalisation).} On the Euclidean section $1/s$ is locally integrable and harmonic off the origin. By the divergence theorem on a ball of radius $r$, $\int\partial^{2}(1/s)\,\varphi \to \varphi(0)\cdot\frac{d}{dr}(r^{-2})\cdot|S^{3}|r^{3} = \varphi(0)(-2r^{-3})(2\pi^{2}r^{3}) = -4\pi^{2}\varphi(0)$, so $\partial^{2}(1/s) = -4\pi^{2}\delta$.

\emph{Step 4 (the vacuum polarisation away from the origin).} The current two-point kernel of the free Dirac field with the coupling of Proposition \claimref{field-theory.coupled-system-tree-level} is $\Pi_{\mu\nu}(x) = -e^{2}\,\mathrm{tr}(\gamma_{\mu}S(x)\gamma_{\nu}S(-x))$ (the minus sign is the fermion loop). With $S(\pm x) = \pm\alpha x + \beta + R(\pm x)$ the terms with an odd number of vector factors have zero trace (Step 1), so
\[
\Pi_{\mu\nu} = 4e^{2}\alpha^{2}(2x_{\mu}x_{\nu} - g_{\mu\nu}s) - 4e^{2}\beta^{2}g_{\mu\nu} + (\text{terms with } R).
\]
At order $m^{0}$ this is $\frac{e^{2}}{\pi^{4}}\frac{2x_{\mu}x_{\nu} - g_{\mu\nu}s}{s^{4}}$, which equals $\frac{e^{2}}{12\pi^{4}}(\partial_{\mu}\partial_{\nu} - g_{\mu\nu}\partial^{2})s^{-2}$ (computed: $(\partial\partial - g\partial^{2})s^{-2} = 12(2xx - gs)/s^{4}$). At order $m^{2}$ it is $\frac{e^{2}m^{2}}{\pi^{4}}(\frac{g}{4s^{2}} - \frac{xx}{s^{3}}) = -\frac{e^{2}m^{2}}{8\pi^{4}}(\partial_{\mu}\partial_{\nu} - g_{\mu\nu}\partial^{2})\frac1s$ (computed). The terms with $R$ and the order-$m^{4}$ products have scaling degree at most $2$ (an $\alpha$ of degree $3$ times the vector part of $R$, degree at most $-1$; $\beta$ of degree $2$ times the scalar part of $R$, degree at most $0$; $m^{4}$ times degree $2$). Both displayed parts are divergence-free (computed), so the remainder $\Pi^{(\le 2)}$ is too, the whole kernel being conserved away from the origin.

\emph{Step 5 (the extension).} Off the origin $\partial^{2}(L_{\mu}/s) = -4/s^{2}$ (computed), so $s^{-2} = -\frac14\partial^{2}(L_{\mu}/s)$, and $L_{\mu}/s$ is locally integrable (the identity of differential renormalisation \cite{FreedmanJohnsonLatorre1992}). Define $\bar\Pi$ as in the statement, all derivatives distributional. On $x \neq 0$ it equals $\Pi$, so it is an extension; it is built from $Q(x)$ and derivatives only, so it is Lorentz covariant (rotation covariant on the Euclidean section, and its Feynman continuation is Lorentz covariant); and $\partial^{\mu}(\partial_{\mu}\partial_{\nu} - g_{\mu\nu}\partial^{2}) = 0$ identically, so the two displayed parts are transverse as distributions. The remainder has scaling degree at most $2 < 4$, so its extension is unique by the imported theorem \cite{BrunettiFredenhagen2000} (Fact \claimref{synthesis.imported-epstein-glaser-extension-theorem-and-ward-preserving-time-ordered-products}(ii)); its divergence is the unique extension of the zero kernel (scaling degree at most $3 < 4$), hence zero. So $\bar\Pi$ is covariant and transverse: the Ward identity holds. The parameter $\mu$ enters only through $L_{\mu}$, and $\mu\,\partial_{\mu}(L_{\mu}/s) = 2/s$, so by Step 3
\[
\mu\,\partial_{\mu}\bar\Pi_{\mu\nu} = -\frac{e^{2}}{48\pi^{4}}(\partial_{\mu}\partial_{\nu} - g_{\mu\nu}\partial^{2})\,2\,\partial^{2}\frac1s = \frac{e^{2}}{6\pi^{2}}(\partial_{\mu}\partial_{\nu} - g_{\mu\nu}\partial^{2})\delta .
\]
Under the Fourier transform $\partial_{\mu} \mapsto I_{4}q_{\mu}$ (on components, numbers read in $\mathrm{span}\{1, I_{4}\}$) the right side is $\frac{e^{2}}{6\pi^{2}}(Q(q)g_{\mu\nu} - q_{\mu}q_{\nu})$; writing $\bar\Pi_{\mu\nu}(q) = (Q(q)g_{\mu\nu} - q_{\mu}q_{\nu})\Pi(q^{2})$, $\mu\,\partial_{\mu}\Pi = e^{2}/(6\pi^{2})$, and since $\Pi$ depends on $\mu$ and $q$ through $\mu^{2}/Q(q)$ alone at leading order, $\Pi = \pm\frac{e^{2}}{12\pi^{2}}\ln(\mu^{2}/Q(q)) + \text{const}$, the sign fixed by the convention relating $\Pi$ to the effective charge (with the fermion-loop sign of Step 4 and $\partial_{\mu} \mapsto I_{4}q_{\mu}$ it is $+$). The magnitude is the coefficient of the stored one-loop result, obtained here with no regulator and with no photon line in the diagram, so neither the Pauli-Villars fields nor the indefinite form of the gauge field enters. The $m^{2}$ part is independent of $\mu$; any other extension of it differs by $c\,g_{\mu\nu}\delta$ (the only covariant order-$0$ ambiguity of a symmetric two-tensor), whose divergence $c\,\partial_{\nu}\delta$ vanishes only for $c = 0$: the transverse extension is unique and contains no photon mass term.

\emph{Step 6 (the ambiguity spaces).} By the imported theorem two extensions of a kernel with degree of divergence $\omega$ differ by $\sum_{|a| \le \omega}c_{a}\partial^{a}\delta$; in momentum space, by a polynomial of degree at most $\omega$ in the momenta, with the kernel's tensor or spinor type. A covariant extension procedure leaves only covariant differences. The check enumerates, for each kernel, the polynomial maps of the right degree and type as linear combinations of monomials in the components of the momenta and polarisations times basis blades of $\Cl(1,3)$, imposes infinitesimal covariance under the six bivector generators $B$ (vectors $v \mapsto \frac12(Bv - vB)$, values $X \mapsto \frac12(BX - XB)$), and computes the dimension of the solution space as the number of unknowns minus an exact rank. The results are those stated in (iii): $3$ for the photon two-point ($g$, $Q(q)g$, $qq$), of which the transverse condition $\Pi(q)(q, b) = 0$ leaves $1$, the combination $Q(q)g - qq$, which is $Z_{3}$; $4$ for the self-energy ($1$, $I_{4}$, $q$, $qI_{4}$, left multiplications), of which the reflection $v \mapsto \gamma_{0}v\gamma_{0}$, $X \mapsto \gamma_{0}X\gamma_{0}$ keeps $1$ and $q$ ($\gamma_{0}I_{4}\gamma_{0} = -I_{4}$), the mass shift $\delta m$ and $Z_{2}$; $2$ for the vertex ($a$, $aI_{4}$), $1$ under the reflection. The Ward-Takahashi identity $q_{\mu}\Lambda^{\mu}(p + q, p) = \Sigma(p) - \Sigma(p + q)$ applied to the ambiguities $\Lambda = c_{V}a$ and $\Sigma = c_{m} + c_{2}q$ reads $c_{V}q = -c_{2}q$: one condition, leaving the two-dimensional space $(c_{m}, c_{2})$ with $c_{V} = -c_{2}$, which is $Z_{1} = Z_{2}$ in the convention $\Lambda_{\mathrm{ct}} = (Z_{1} - 1)\gamma$, $\Sigma_{\mathrm{ct}} \ni -(Z_{2} - 1)q$ (computed jointly: dimension $2$ under the reflection, $4$ without, the two pseudoscalar terms being tied in the same way). For three photons the covariant trilinear forms of degree at most $1$ in $(q_{1}, q_{2})$, $q_{3} = -q_{1} - q_{2}$, span $8$ dimensions ($(a_{i}\cdot a_{j})(a_{k}\cdot q)$ and $\epsilon(a_{1}, a_{2}, a_{3}, q)$, $q \in \{q_{1}, q_{2}\}$), and the transverse conditions on any two legs ($a_{i} = q_{i}$) leave none; for four photons the constant covariant forms span $4$ dimensions (three products of $g$ and one $\epsilon$), and transversality on one leg leaves none. Hence, if covariant extensions satisfying the Ward identities and the reflection exist, any two differ exactly by $Z_{3}$, $\delta m$ and $Z_{2} = Z_{1}$, the counterterms of electrodynamics, and Furry's theorem is not needed to fix the three-photon kernel. Existence is proved here only for the vacuum polarisation (Step 5); for the self-energy and vertex, the four-photon and the three-current kernels it is the known result of the causal approach (Scharf \cite{Scharf1995}; Fact \claimref{synthesis.imported-epstein-glaser-extension-theorem-and-ward-preserving-time-ordered-products}(iii)), not proved here. The self-energy and vertex kernels contain the photon propagator in a covariant gauge, so these parts are read under the Krein-form premise \claimref{synthesis.krein-form-premise-indefinite-inner-product-for-gauge-fields}; the counting itself uses only covariance. The same eight-dimensional space contains every covariant ambiguity of a three-current kernel with one axial and two vector legs (the $\epsilon$ terms included), so at most one covariant extension conserves the two vector currents, and every such extension (one exists by the imported result) has the same axial divergence; the value computed with Pauli-Villars fields in the triangle result \claimref{synthesis.triangle-anomaly-in-the-spacetime-algebra-is-the-pseudoscalar-part-of-the-field-square-and-cancels-per-generation} is therefore that of every covariant vector-conserving scheme.

\emph{Step 7 (separately, the scalar bubble of a free scalar field).} $\Delta_{0}^{2} = 1/(16\pi^{4}s^{2}) = -\frac{1}{64\pi^{4}}\partial^{2}(L_{\mu}/s)$ off the origin, so its extension has $\mu\,\partial_{\mu} = -\frac{1}{64\pi^{4}}\cdot 2\cdot(-4\pi^{2})\delta = \delta/(8\pi^{2})$. For $\lambda\phi^{4}/4!$ the one-loop four-point function is $\lambda - \frac{\lambda^{2}}{2}\sum_{\text{3 channels}}\bar B$, and its independence of $\mu$ at this order gives $\mu\,d\lambda/d\mu = 3\cdot\frac12\cdot\frac{1}{8\pi^{2}}\lambda^{2} = \frac{3\lambda^{2}}{16\pi^{2}}$.

\emph{Step 8 (merged).} Clause (vii) points to Theorem \claimref{synthesis.floor-regulator-keeps-the-charge-and-the-central-term-below-the-cut-off}(v), into which the statement on the floors was merged on 2 November 2026; its closure and Schur steps, with the irreducibility of the attachment's modules (Proposition \claimref{field-theory.twistor-attachment}), are proved there. The check's function for it is kept as a cross-check: it recomputes the seven-dimensional commutant of $H = \gamma_{0}e_{+}$ and its closure with seven sample generators.

\emph{Step 9 (the units).} The conventions place the field theory's unit outside the Clifford product on massive lines (the right unit), which commutes with every left multiplication. The ambiguities of Step 6 are left multiplications, so the terms $I_{4}$ and $qI_{4}$ are pseudoscalar terms (the matrix reading $\gamma_{5}$), odd under the reflection; complex rescalings of $\delta m$ and $Z_{2}$ by the right unit are a separate freedom, removed by the reality of the coefficients. Had the unit been placed as $I_{4}$ inside the product, the two would not be distinguishable on odd elements.

\emph{Step 10 (on a gravitational field; argued).} Theorem \claimref{synthesis.hadamard-dirac-states-scale-to-the-flat-massless-vacuum-at-every-point-of-a-gravitational-field} gives, at every point of a globally hyperbolic gravitational field, $l^{3}$ times the Dirac two-point function at scale $l$ tending to the flat massless kernel: scaling degree $3$ with the flat limit. The time-ordered kernels built from Hadamard parametrices are defined off the diagonal and their products obey the scaling-degree bound by the wave front set condition \cite{BrunettiFredenhagen2000}; the leading part of each one-loop kernel is then the flat one, so its degree of divergence and its leading scale coefficient are the flat values, and the remaining freedom adds local terms in the field strength of matching dimension, as in the curved-field construction of Hollands and Wald \cite{HollandsWald2001} (cited from memory). This step is argued, not proved.
\end{proof}
